\section{Full robustness tables}
\label[appendix]{app:robustness}

These tables report every robustness check with the three inferences. Only the bootstrap by year is used for conclusions; the other two are shown for comparison, since the size simulations of \cref{sec:inference} show that they over-reject on these data.

\begin{table}[H]
\centering
\footnotesize
\caption{Wheat 2010--2025 (phase 4): heat coefficient in every specification}
\label{tab:rob-wheat-full}
\begin{threeparttable}
{\scriptsize\setlength{\tabcolsep}{3.5pt}% GÉNÉRÉ par scripts/26_tables_memoire.py (phase4_robustesse.csv)
\begin{tabular}{llllrrr}
\toprule
& Variant & Term & $\hat\beta$ & $p$ clust.\ dept. & $p$ Conley & $p$ boot.\ (two-sided) \\
\midrule
 & M1 (reference) & DD$_{>32}$, May--June & $-0.0149$ & 0.004 & 0.047 & 0.374 \\
R1 & Unweighted & DD$_{>32}$, May--June & $-0.0159$ & 0.034 & --- & 0.332 \\
R2 & Without 84, 30, 48 & DD$_{>32}$, May--June & $-0.0147$ & 0.005 & 0.055 & 0.385 \\
R3 & Without 2025 & DD$_{>32}$, May--June & $-0.0290$ & <0.001 & <0.001 & 0.141 \\
R4 & 85 departments (status 1 and 2) & DD$_{>32}$, May--June & $-0.0151$ & 0.003 & 0.042 & 0.367 \\
R5 & Threshold 30\,$^{\circ}$C & DD$_{>30}$, May--June & $-0.0070$ & 0.004 & 0.076 & 0.398 \\
R6 & Department-area weather & DD$_{>32}$, May--June & $-0.0127$ & 0.006 & 0.036 & 0.442 \\
R7 & + regional linear trends & DD$_{>32}$, May--June & $-0.0125$ & 0.038 & --- & 0.444 \\
R8 & Long panel 1990--2025 & DD$_{>32}$, May--June & $-0.0100$ & 0.038 & 0.171 & 0.416 \\
R9 & Conley 300 km & DD$_{>32}$, May--June & $-0.0149$ & 0.004 & --- & 0.374 \\
R10 & Window May--July & DD$_{>32}$, May--July & $+0.0036$ & 0.243 & 0.499 & 0.663 \\
R11 & Window shifted $-2$ weeks & DD$_{>32}$, window $-2$ weeks & $-0.1582$ & <0.001 & 0.014 & 0.342 \\
R12 & Window shifted $+2$ weeks & DD$_{>32}$, window $+2$ weeks & $-0.0065$ & 0.167 & 0.354 & 0.573 \\
R13 & Without mountain departments & DD$_{>32}$, May--June & $-0.0152$ & 0.004 & 0.050 & 0.377 \\
R14 & Temperature bins & Days 30--33\,$^{\circ}$C & $-0.0170$ & 0.291 & --- & 0.738 \\
R14 & Temperature bins & Days $\geq 33$\,$^{\circ}$C & $-0.0206$ & 0.298 & --- & 0.697 \\
\bottomrule
\end{tabular}
}
\begin{tablenotes}\footnotesize
\item R1--R10 were pre-specified; R11--R14 were added after the first estimates and defined before being computed. R2 drops Vaucluse (84), Gard (30) and Lozère (48), small departments that carry much of the unweighted variation in heat. R9 changes only the Conley cut-off. ``---'': not reported, because \texttt{fixest} had to correct a Conley matrix that was not positive definite. R14 reports two temperature bins, 30--33\,$^{\circ}$C and above 33\,$^{\circ}$C, in days.
\end{tablenotes}
\end{threeparttable}
\end{table}

\begin{table}[H]
\centering
\footnotesize
\caption{Maize 2000--2025 (phase 4 bis): main coefficient in every specification}
\label{tab:rob-maize-full}
\begin{threeparttable}
% GÉNÉRÉ par scripts/26_tables_memoire.py
\begin{tabular}{lrrrrr}
\toprule
Variant & $\hat\beta$ & $p$ clust.\ dept. & $p$ Conley & $p$ boot.\ two-sided & $p$ boot.\ one-sided \\
\midrule
Main & $+0.0158$ & 0.009 & --- & 0.016 & 0.007 \\
38 departments & $+0.0217$ & <0.001 & <0.001 & 0.010 & 0.006 \\
2010--2025 & $+0.0242$ & 0.006 & --- & 0.056 & 0.030 \\
Department-area weather & $+0.0104$ & 0.024 & <0.001 & 0.441 & 0.221 \\
Threshold 30\,$^{\circ}$C & $+0.0104$ & 0.004 & --- & 0.016 & 0.007 \\
July only & $+0.0575$ & 0.001 & <0.001 & 0.029 & 0.013 \\
Unweighted & $+0.0185$ & 0.008 & <0.001 & 0.045 & 0.021 \\
Without 2025 & $+0.0157$ & 0.018 & --- & 0.031 & 0.015 \\
Without 2003 & $+0.0243$ & 0.008 & <0.001 & 0.060 & 0.029 \\
2008 mask (blocks) & $+0.0153$ & 0.012 & --- & 0.015 & 0.007 \\
\bottomrule
\end{tabular}

\begin{tablenotes}\footnotesize
\item One-sided: in the direction of the hypothesis $\beta > 0$. ``---'': Conley matrix not positive definite.
\end{tablenotes}
\end{threeparttable}
\end{table}
